\documentclass[11pt,a4paper]{article}
\usepackage[T1]{fontenc}
\usepackage[utf8]{inputenc}
\usepackage{amsmath,amssymb,amsthm}
\usepackage[margin=1in]{geometry}
\usepackage[colorlinks=true,linkcolor=blue,urlcolor=blue]{hyperref}
\theoremstyle{plain}
\newtheorem{theorem}{Theorem}
\newtheorem{lemma}[theorem]{Lemma}
\newtheorem{proposition}[theorem]{Proposition}
\newtheorem{corollary}[theorem]{Corollary}
\theoremstyle{definition}
\newtheorem{remark}[theorem]{Remark}

\title{The empirical recurrence for OEIS A232033, proved by transfer matrix}
\author{Adrian Perez Fontelles\\ \small Independent researcher}
\date{3 September 2026}
\begin{document}
\maketitle

\begin{abstract}
OEIS A232033 counts arrays of width $4$ over $\{0,\dots,1\}$ under a condition on each cell
and its neighbours, and carries an empirical recurrence of order $20$ contributed by
R.~H.~Hardin. It is true, and it is decidable rather than empirical. Every neighbour named lies
in the three rows $i-1,i,i+1$, so an ordered pair of consecutive rows is a state, the count is
a walk count on $S=272$ of them, and the conjectured recurrence is then settled by a finite
exact computation.
\end{abstract}

\noindent\small 2020 Mathematics Subject Classification. 05A15, 05B45, 68Q45.\normalsize

\section{The conjecture}

OEIS A232033 is ``Number of (n+1)X(3+1) 0..1 arrays with every element equal to some horizontal, vertical or antidiagonal neighbor, with top left element zero.''. It has offset $1$ and begins
\[
37, 521, 6958, 94102, 1272579, 17207275, 232668433, 3146033503,\ \dots
\]
The entry states, as a conjecture contributed by its author and never marked settled:
\begin{quote}
Empirical: a(n) = 16*a(n-1) -32*a(n-2) -23*a(n-3) +30*a(n-4) +49*a(n-5) +160*a(n-6) -163*a(n-7) -304*a(n-8) +146*a(n-9) +59*a(n-10) +316*a(n-11) -147*a(n-12) -255*a(n-13) +70*a(n-14) +122*a(n-15) -21*a(n-16) -28*a(n-17) +2*a(n-18) +5*a(n-19) -a(n-20)
\end{quote}
As of the ``Last modified'' line on the live entry (Jul 23 2025 06:49:05, revision 6) this is still
recorded as empirical, and nothing on the entry records it as proved.

\section{What is being counted}

The neighbours of a cell are the cells at the offsets
\[
(-1,0), (-1,1), (0,-1), (0,1), (1,-1), (1,0),
\]
those falling outside the array being absent. The entry requires that every cell equals at least one of its neighbours. The entry also fixes the top left cell to zero, which restricts the states an array may begin in and nothing else.

\section{The count is a walk count}

Every offset above moves the row index by at most one, so whether a cell of row $i$ passes
depends on rows $i-1$, $i$ and $i+1$ and nothing else. Take as vertices the ordered PAIRS
$(p,c)$ of consecutive rows, allowing $p$ to be ABSENT --- that is what lets a walk of no steps
describe a one-row array --- and put an edge $(p,c)\to(c,x)$ exactly when every cell of the
middle row $c$ passes with $p$ above it and $x$ below it. A state with $p$ absent may begin an
array; a state $(p,c)$ accepts when every cell of $c$ passes with nothing below it. An array of
$L$ rows is then a walk of $L-1$ steps, each step settling one row, and every array arises from
exactly one walk. Hence
\[
a(n)\;=\;\iota^{\!\top}M^{\,j}\tau
\]
with $\iota$ the indicator of the beginning states, $\tau$ of the accepting ones, and $j$ the
walk index, which differs from the entry's own $n$ by a constant read off from its published
terms. There are $S=272$ states, and $a$ is $C$-finite.

At $n=1$ the array has one row, and
the condition is tested there with nothing above or below it; the entry's first term is
$37$, which the model reproduces along with every later term. That is the cheapest check
that the neighbour set and the boundary have been read the same way the entry reads them: a
condition tested as though the missing rows were present, rather than absent, gives a different
first term. There are $2^{4}=16$ possible rows over the entry's own alphabet. The state count $S=272$ is exactly $16\cdot17$,
one state for each ordered pair of a row and its possible predecessor, the absent predecessor
included.

\section{The proof}

\begin{lemma}\label{lem:crit}
Let $q(t)=t^{20}-\sum_i c_it^{20-i}$ be the characteristic polynomial of the
conjectured recurrence. The residual $a(n)-\sum_ic_ia(n-i)$ equals
$\iota^{\!\top}M^{\,j}q(M)\tau$ for the corresponding walk index $j$, so the recurrence
holds from some point on if and only if $\iota^{\!\top}M^{j}q(M)\tau=0$ for all large $j$,
and it is enough to examine $j<S$.
\end{lemma}

\begin{proof}
Factor $M^{j}$ out of $M^{j+20}-\sum_ic_iM^{j+20-i}$. For the bound, $M^{S}$
is an integer combination of $I,M,\dots,M^{S-1}$ by Cayley--Hamilton, so the numbers
$u_j=\iota^{\!\top}M^{j}q(M)\tau$ obey the monic recurrence given by the characteristic
polynomial of $M$; $S$ consecutive zeros force every later one, and the last nonzero $u_j$
pins the threshold exactly.
\end{proof}

\begin{theorem}
\[
a(n)\;=\;16\,a(n-1) - 32\,a(n-2) - 23\,a(n-3) + 30\,a(n-4) + 49\,a(n-5) + 160\,a(n-6) - 163\,a(n-7) - 304\,a(n-8) + 146\,a(n-9) + 59\,a(n-10) + 316\,a(n-11) - 147\,a(n-12) - 255\,a(n-13) + 70\,a(n-14) + 122\,a(n-15) - 21\,a(n-16) - 28\,a(n-17) + 2\,a(n-18) + 5\,a(n-19) - a(n-20)
\]
for every $n>20$.
\end{theorem}

\begin{proof}
The vector $w=q(M)\tau$ was formed by $20$ matrix--vector products in exact integer
arithmetic and $\iota^{\!\top}M^{j}w$ evaluated for $j=0,1,\dots$, again exactly; those
integers vanish from the index corresponding to $n=20+1$ onwards and the run continued
until the working vector was identically zero, which settles every larger $j$ at once.
\end{proof}

\section{Verification}

Three checks, each independent of the algebra above.

First, the model reproduces all $16$ terms the entry publishes, exactly, in integer
arithmetic. This is what ties the model to the entry: the digraph is built by reading the
entry's English, and a misreading gives different counts.

Second, the conjectured recurrence was evaluated directly on those published terms, with no
matrices involved, and holds wherever the proved range and the data overlap.

Third, the annihilation test of Lemma~\ref{lem:crit} was carried out in exact integer
arithmetic on the whole vector rather than on a sampled prefix.

\begin{thebibliography}{9}
\bibitem{oeis} The OEIS Foundation, \emph{The On-Line Encyclopedia of Integer
Sequences}, \url{https://oeis.org}, 2026. Sequence A232033.
\bibitem{stanley} R.~P.~Stanley, \emph{Enumerative Combinatorics}, Volume 1, 2nd ed.,
Cambridge University Press, 2011. (Transfer-matrix method, Section 4.7.)
\bibitem{comtet} L.~Comtet, \emph{Advanced Combinatorics}, Reidel, 1974.
\bibitem{horn} R.~A.~Horn and C.~R.~Johnson, \emph{Matrix Analysis}, 2nd ed., Cambridge
University Press, 2013.
\end{thebibliography}

\end{document}
